% CVPR 2024 style report — What Happens Next? (Track A + Track B)
% Compile : pdflatex report.tex  (run from the repo root so figures/ resolves)

\documentclass[10pt,twocolumn,letterpaper]{article}

\usepackage{cvpr}              % style file from CVPR template
\usepackage{times}
\usepackage{epsfig}
\usepackage{graphicx}
\usepackage{amsmath,amssymb}
\usepackage{booktabs}
\usepackage{xcolor}
\usepackage{float}
\usepackage[breaklinks=true,bookmarks=false]{hyperref}

\cvprfinalcopy   % retirer pour la version anonyme
\def\cvprPaperID{****}
\def\httilde{\mbox{\tt\raisebox{-.5ex}{\symbol{126}}}}

\ifcvprfinal\pagestyle{empty}\fi

% Graphics path so \includegraphics{x.png} finds files alongside report.tex
\graphicspath{{./}{figures/}}

% Compile-safe figure include: shows a placeholder box if the file is absent,
% so the document always builds even before every figure has been generated.
% Probes both ./ and figures/ (graphicspath then resolves the actual include).
\newcommand{\safefig}[2][\columnwidth]{%
  \IfFileExists{#2}%
    {\includegraphics[width=#1]{#2}}%
    {\IfFileExists{figures/#2}%
      {\includegraphics[width=#1]{#2}}%
      {\fbox{\begin{minipage}[c][3.0cm][c]{#1}\centering\ttfamily\scriptsize
         [figure not found]\\\detokenize{#2}\end{minipage}}}}%
}

\begin{document}

%==================================================================
\title{Action recognition on a Curated Something-Something\,V2 Subset:\\ What Happens Next?}

\author{Etienne CHEVROLAT \quad Alfred RUSCHER\\
\'Ecole polytechnique}

\maketitle

%==================================================================
\begin{abstract}
We study the two subproblems of the \emph{What Happens Next} challenge:
classify human-object interactions into 33 SSV2-derived classes from
clips of \textbf{only 4 frames} per video. In \textbf{Track A} (no external
data) we design and train four architectures entirely from scratch --- a
CNN$+$Transformer family, an R(2+1)D backbone coupled to a spatio-temporal
Transformer, and a vanilla ViT-B initialised by a \emph{from-scratch}
VideoMAE\,v2 pre-training on the challenge data --- and show that held-out
accuracy saturates near~$0.41$: the R(2+1)D$+$Transformer reaches~$0.406$
at $39.8$\,M parameters and $4.2$\,ms/clip, on par with the $86$\,M
self-supervised ViT-B ($0.409$). In \textbf{Track B} we benchmark 11
architectures of widely different sizes ($36$\,M to $7$\,B parameters) and
five families of overfitting mitigations on top of a VideoMAE-B baseline
that reaches~$0.6096$. Our best single model is a head-only fine-tune of
V-JEPA\,2 ViT-L SSV2 with linear temporal upsampling, reaching~$0.638$. The
$\sim$20\,pt gap between the best from-scratch and the best pre-trained
model isolates \emph{pre-training, not architecture}, as the dominant factor
at this data and frame budget.
\end{abstract}

%==================================================================
\section{Problem statement}
\label{sec:intro}

The dataset is a curated 33-class subset of Something-Something\,V2
(SSV2)~\cite{goyal2017something}, organised as $44{,}993$ training
clips (none of which are labelled 27), $6{,}746$ validation (labelled) clips and $6{,}913$ test clips (unlabelled) to create a graded submission for the challenge. Each clip is delivered as
4 ($224\times 224$) JPEG frames, which is far below the temporal
resolution at which video foundation models are typically
pre-trained~(16--64 frames) as well as below their spatial resolution training (below their training spatial resolution ($384\times 384$ for V-JEPA 2 ViT-g fpc64-384, for instance)). The challenge therefore stresses
\emph{(i)}~transfer from a frame-rich architecture to a sparse
deployment regime, and \emph{(ii)}~generalisation under a small
labelled-data budget.

%==================================================================
\section{Method}
\label{sec:method}

\subsection{Track B}

\subsubsection{Architectures surveyed}
\label{sec:method:archs}
VideoMAE~\cite{tong2022videomae} specialises in reconstructing
hidden pixels. As our problem is characterised by the lack of
frames, we first focused on the 86\,M-parameter VideoMAE-B, before
trying bigger models and finally focusing on
V-JEPA\,2~\cite{bardes2024vjepa2}.

All in all, we tested 11 families spanning four design philosophies:
\emph{(a)~masked-pixel pretraining} (VideoMAE-B/L/v2-G~\cite{tong2022videomae},
Hiera-B/Huge~\cite{ryali2023hiera}); \emph{(b)~joint-embedding predictive architectures}
(V-JEPA\,2 ViT-L and ViT-g~\cite{bardes2024vjepa2}); \emph{(c)~image-language contrastive
backbones} (X-CLIP~\cite{ni2022xclip}, SigLIP-2 Giant~\cite{tschannen2025siglip2},
DINOv3~\cite{oquab2023dinov2}); and \emph{(d)~attention-only video transformers}
(TimeSformer~\cite{bertasius2021timesformer}, MViT-v2~\cite{fan2021multiscale},
Swin3D~\cite{liu2022videoswin}, InternVideo2~6B~\cite{wang2024internvideo2}). We also
evaluated two vision-language models, Qwen-VL~7B~\cite{bai2023qwenvl} and
LLaVA-Mistral~7B~\cite{liu2023llava}. For each family we kept the public
SSV2-finetuned weights when available. For backbones without an SSV2 checkpoint we used
the closest available (Kinetics-400 or DINO/SigLIP raw features) and
added the temporal pooler at the head level. The results are visible in Figure~\ref{fig:acc-vs-params}.

\subsubsection{Anti-overfit toolbox}
\label{sec:method:overfit}
After optimizing the hyperparameters of VideoMAE-B, we encountered a plateau where the experiments would reach a training accuracy higher than $0.70$ in 2-3 epochs, while the validation accuracy would remain below $0.61$. We tried different techniques to bridge this generalization gap:

\paragraph{Classical regularisation} Dropout
($\{0.05,...,0.50\}$), weight decay
($\{10^{-5},...,10^{-2}\}$), label smoothing
($\{0,...,0.50\}$) and gradient clipping
($\|g\|_2 \le 1$); see Sec.~\ref{sec:results}.

\paragraph{Parameter-efficient fine-tuning}
\emph{(i)}~LoRA adapters of rank
$r \in \{8,16,32,64,128\}$ on all 40 blocks of V-JEPA\,2 ViT-g (180\,M
trainable for $r{=}16$); \emph{(ii)}~block freezing controlled by
\textit{num\_frozen\_blocks} $\in \{0,...,40\}$; \emph{(iii)}~layer-wise
LR decay (LLRD) with factor~$0.80$; \emph{(iv)}~differential learning
rates between the backbone and the new head, parameterised by
\textit{backbone\_lr\_scale} $\in [0.02,1.0]$; \emph{(v)} Head MLP / attention pool / cross attention pair; \emph{(vi)} Multi-scale temporal features and number of temporal layers; \emph{(vii)} Choice of optimizer (AdamW, Prodigy, gradient accumulation, etc.)

\paragraph{Augmentation} Random resized crop, horizontal flip
($p{=}0.5$, not to be applied to classes 18 "pulling something from left to right" and 19 "pulling something from right to left"), colour jitter ($p{=}0.8$, strength $0.35$), Gaussian blur
($p{=}0.1$), and temporal equivariant reverse, where reversing
the clip flips the label between matched action pairs (e.g. "opening something"
(class 10) $\leftrightarrow$ "closing something"
(class 0)). For input mixing we use
\textbf{MixUp}~\cite{zhang2018mixup}:
\begin{equation}
\tilde{x} = \lambda x_i + (1-\lambda) x_j,\quad
\tilde{y} = \lambda y_i + (1-\lambda) y_j,
\end{equation}
with $\lambda \sim \text{Beta}(\alpha,\alpha)$, $\alpha=0.4$.

\paragraph{Trajectory stabilisation and Class re-weighting} EMA on the model weights with
decay~$0.9998$, cosine LR schedule with $3$ warm-up epochs and a
floor of~$10^{-7}$. Inverse-frequency class weights to
counter the long-tail distribution of Track B.

\subsection{Going beyond VideoMAE-B}
\label{sec:method:beyond}
Following the intermediary results presentation by the challenge's best teams, two main directions were explored:
\emph{(a)}~LoRA on V-JEPA\,2 ViT-g with four head/LoRA learning-rate
combinations (a~$3\times 3$ grid over $\{10^{-3},10^{-2}\}$ for the head
and a multiplicative scale~$\in[0.05,1]$ for LoRA); and
\emph{(b)}~head-only fine-tuning of V-JEPA\,2 ViT-L SSV2 where the 4
input frames are linearly interpolated (LERP) to fill the~16 input
positions expected by the tube tokeniser, restricting the
174-way head to the 33 SSV2 templates that appear in Track B.

\subsection{Track A: learning from scratch}
\label{sec:method:trackA}

Track A forbids external pre-training, so every weight is learned on the
$44{,}993$ labelled clips of the curated SSV2 subset, at the same 4-frame
budget as Track B. We explored a deliberate capacity / inductive-bias
ladder of four custom architectures (all implemented from scratch in
\texttt{src/models/}).

\paragraph{(1) CNN spatial encoder $\to$ temporal Transformer (12.9\,M).}
A shared ResNet-18-style encoder (7$\times$7 stem, four 2-block stages,
$64\!\to\!512$ channels) maps each of the 4 frames to a single vector via
global average pooling. The 4 vectors receive a learned temporal position
embedding and pass through 2 self-attention blocks; mean-pooling feeds a
linear head. The temporal Transformer sees only $T{=}4$ tokens, so its
capacity is largely under-used.

\paragraph{(2) CNN spatial encoder $\to$ spatio-temporal Transformer (17.8\,M).}
We drop the spatial pooling and keep the $7\times7$ feature grid, yielding
$4\times49 = 196$ spatio-temporal tokens with a \emph{factorised} position
embedding $\mathrm{PE}(t,p)=\mathrm{pos}_t+\mathrm{pos}_p$. A joint
Transformer then reasons over space and time. Despite $16\times$ more
tokens, accuracy does \emph{not} improve: a ResNet-18 trained from scratch
on $45$k clips is data-starved, and extra tokens cannot manufacture the
motion signal the convolutions failed to extract.

\paragraph{(3) R(2+1)D backbone + spatio-temporal Transformer (39.8\,M).}
We replace the 2D CNN with an R(2+1)D-18 stem~\cite{tran2018closer}, which
factorises every 3D convolution into a spatial $(1,3,3)$ and a temporal
$(3,1,1)$ kernel and therefore models motion natively. Three targeted
additions bridge the small-$T$ regime: \emph{(i)}~a parameter-free Temporal
Shift Module~\cite{lin2019tsm} in every residual block, forcing early
cross-frame mixing; \emph{(ii)}~removal of the spatial pool so the stage-4
output is flattened into $4\times7\times7=196$ tokens fed to a 3-layer
factorised-PE Transformer; \emph{(iii)}~a learned attention-pooling head in
place of a CLS token. This model is initialised from the plain R(2+1)D run
and trained with a 5-epoch backbone freeze, MixUp/CutMix, stochastic depth
and \texttt{Dropout3d}.

\paragraph{(4) Self-supervised MAE pre-training + vanilla ViT-B (86.3\,M).}
Finally we test whether self-supervision \emph{on the challenge data alone}
can rival external pre-training. We re-implement VideoMAE\,v2~\cite{tong2022videomae}
from scratch --- tube masking, dual masking, a light 4-block decoder and a
norm-pixel reconstruction loss --- and pre-train its ViT-B encoder
(tubelet $2\times16\times16$, $\dim{=}768$, depth 12) on the
\emph{unlabelled} training frames. The encoder is then loaded $100\%$ into a
ViT-B classifier and fine-tuned with layer-wise LR decay (LLRD~$0.75$),
5-epoch warm-up, cosine schedule, MixUp/CutMix and label smoothing $0.1$.

%==================================================================
\section{Results}
\label{sec:results}

\subsection{Scaling laws on Track B}

\begin{figure}[t]
  \centering
  \safefig{fig_acc_vs_params.png}
  \caption{Top-1 \emph{val\_dir} accuracy as a function of backbone size,
           best result per family. Small dots show individual runs;
           large markers indicate the family best. The dashed
           line~$0.638$ corresponds to the V-JEPA\,2 ViT-L
           single-model record (305\,M parameters).}
  \label{fig:acc-vs-params}
\end{figure}

Figure~\ref{fig:acc-vs-params} establishes the headline finding of
our study: on Track B, accuracy saturates around 300--600\,M
parameters and then plateaus, even though some pre-trained models
reach 6--7\,B parameters. The image-language contrastive backbones
(SigLIP-2, DINOv3, X-CLIP) cluster between 0.46 and 0.53. The
SSV2-pretrained masked-pixel models reach 0.59--0.61. V-JEPA\,2 ViT-L
is the only family that exceeds the VideoMAE-B baseline, at 305\,M
parameters.

\subsection{Single-axis sweeps on VideoMAE-B}
\begin{figure*}[t]
  \centering
  \safefig[\textwidth]{fig_opt_sweeps.png}
  \caption{Single-axis sweeps starting from the best VideoMAE-B
           configuration (other axes fixed at their optimum). The dashed
           gold line shows the joint optimum~$0.6086$ obtained when all
           eight axes are simultaneously optimised. Error bars are
           $\pm 1\sigma$ over three random seeds.}
  \label{fig:sweeps}
\end{figure*}
Figure~\ref{fig:sweeps} summarises the eight single-axis sweeps.
A few  remarks:
(a)~head width plateaus at $4096$, with a slight \emph{drop} at $8192$
attributable to overfitting on the training set;
(b)~dropout has a flat optimum around $0.10$--$0.15$;
(c)~label smoothing peaks at $0.20$, in line with~\cite{szegedy2016rethinking};
(d)~weight decay shows the classic $5\!\cdot\!10^{-4}$ optimum;
(g)~freezing any of the 12 VideoMAE-B blocks already costs accuracy,
which contrasts with our V-JEPA\,2 experiments where freezing 32 of
40 blocks helps. The capacity-to-data ratio is more sensitive for
VideoMAE-B.

The combination of the optimized hyperparameters gives us the configuration
\begin{center}\small
\begin{tabular}{@{}rl@{\quad}rl@{}}
$h$           & 4096                 & dropout         & 0.10 \\
label smooth. & 0.20                 & weight decay    & $5{\cdot}10^{-4}$ \\
epochs        & 30                   & bb\_lr\_scale   & 0.5 \\
frozen blocks & 0                    & frames          & 4 \\
EMA decay     & 0.9998               &                 & \\
\end{tabular}
\end{center}
yielding $0.6086$ top-1 accuracy.

\subsection{Anatomy of the best run}

\begin{figure}[t]
  \centering
  \safefig{fig_opt_curves.png}
  \caption{Train, internal-val and held-out (\emph{val\_dir}) accuracy
           and loss during the best VideoMAE-B run. The gold star marks
           the early-stopping checkpoint at epoch 30.}
  \label{fig:curves}
\end{figure}

Figure~\ref{fig:curves} shows that the model continues to fit the
training set well past the point at which the held-out accuracy
plateaus, leading to a residual train/val gap of $\sim$28\,pts. Adding
mixup, EMA and stronger label smoothing reduces this gap by
$\sim$3\,pts but does not close it.

\begin{figure}[t]
  \centering
  \safefig{fig_best_confusion.png}
  \caption{Confusion matrix of the best VideoMAE-B configuration on
           \emph{val\_dir}. The six pairs of classes most often confused
           are marked with a red cross; they all correspond to
           visually similar transitive actions (e.g.\ ``moving X''
           vs.\ ``putting X'').}
  \label{fig:confusion}
\end{figure}

Figure~\ref{fig:confusion} illustrates the dominant residual errors.
Five of the 33 classes account for over a third of the remaining
errors; these are the classes whose template description differs by
a single preposition (\emph{onto} vs.\ \emph{into}) or by direction
(\emph{up} vs.\ \emph{down}). These ambiguities are exactly the cases
where having only 4 frames is most penalising.
Another result is that some on the temporally exchanged class (0 and 10 for instance) are both in the most baddly predicted ones. We stopped using this augmentation and the model yielded an accuracy of~$0.6096$

\subsection{Going beyond: LoRA and head-only V-JEPA\,2}

LoRA on V-JEPA\,2 ViT-g converges quickly (validation accuracy of
$0.64$ after a single epoch with head learning rate~$10^{-3}$,
backbone-LR scale~$1.0$), but does not exceed the head-only
V-JEPA\,2 ViT-L recipe (Section~\ref{sec:method:beyond}), which
attains $\mathbf{0.638}$ top-1 on \emph{val\_dir} in two epochs;
the full grid of tested configurations is shown in
Fig.~\ref{fig:lora}.
A frame-interpolation ablation shows that linear interpolation
between the four source frames underperforms against naive frame
repetition ($-6.6$\,pts) and is on par with Gaussian temporal
blending~($\pm 0.5$\,pt) (see Fig.~\ref{fig:upsampling} for the different interpolation tested, with the resulting top1 valdir accuracy of VideoMAE-B best configuration in Table \ref{tab:upsampling} ). This could be because the models were
pretrained on clear-cut pictures. Adding 4-clip averaging during
inference \emph{degrades} accuracy on this dataset, contrary to
typical SSV2 evaluation protocols, presumably because the action of
interest spans the whole clip.

\begin{figure}[!ht]
  \centering
  \safefig{lora_sweep_results.png}
  \caption{LoRA fine-tuning sweep: best \emph{val\_dir} top-1
           accuracy per architecture~$\times$~LoRA rank.}
  \label{fig:lora}
\end{figure}

\subsection{Track A: results}
\label{sec:results:trackA}

\begin{figure}[t]
  \centering
  \safefig{fig_trackA_acc_vs_params.png}
  \caption{Track A held-out top-1 vs.\ trainable parameters (log scale).
           Small dots are individual runs, large markers the family best.
           Dashed lines recall the \emph{pre-trained} Track B references
           (VideoMAE-B $0.6096$, V-JEPA\,2 ViT-L $0.638$). From-scratch
           accuracy saturates near $0.41$, $\sim$20\,pts below any
           pre-trained backbone.}
  \label{fig:trackA-scaling}
\end{figure}

\begin{table}[t]
  \centering\small
  \begin{tabular}{@{}llcccc@{}}
    \toprule
    Model & Architecture & Params & top-1 & top-5 & ms/clip \\
    \midrule
    m1.2 & CNN $\to$ ViT$_t$              & 12.9\,M & 0.244 & 0.548 & 1.3 \\
    m2   & CNN $\to$ ST-ViT               & 17.8\,M & 0.229 & 0.535 & 1.6 \\
    m5B  & R(2+1)D+TSM+ST-ViT             & 39.8\,M & \textbf{0.406} & 0.719 & 4.2 \\
    m7   & MAE-ViT-B (self-sup.\ on train)& 86.3\,M & \textbf{0.409} & 0.736 & 9.2 \\
    \bottomrule
  \end{tabular}
  \caption{From-scratch Track A models on the held-out \emph{val\_dir}
           ($6{,}745$ clips). Best held-out checkpoint per family;
           latency is forward-only per clip on a single GPU.}
  \label{tab:trackA}
\end{table}

Figure~\ref{fig:trackA-scaling} and Table~\ref{tab:trackA} give the
headline result for Track A: \emph{without} external pre-training,
accuracy plateaus around $0.41$ and never approaches the $0.61$ reached
by an SSV2-pre-trained VideoMAE-B. The R(2+1)D$+$Transformer (m5B) and the
self-supervised ViT-B (m7) tie within noise ($0.406$ vs.\ $0.409$), but
m5B reaches it with $2.2\times$ fewer parameters and at less than half the
latency ($4.2$ vs.\ $9.2$\,ms), making it the best accuracy/compute
trade-off of the track. The two small CNN$+$Transformer models
($\le 18$\,M) lag by $16$--$18$\,pts, confirming that capacity \emph{and}
an explicit motion prior are both needed in this regime.

\begin{figure*}[t]
  \centering
  \safefig[\textwidth]{fig_trackA_model7_curves.png}
  \caption{model7\_A fine-tuning. \emph{Left:} layer-wise LR decay is
           decisive --- without it the held-out accuracy plateaus at
           $0.26$, with it (plus diverse augmentation) it reaches $0.41$.
           \emph{Right:} on a clean 50-epoch run, the internal-val curve
           sits $\sim$17\,pts above the held-out curve, the same
           generalisation gap that dominates Track B.}
  \label{fig:trackA-m7}
\end{figure*}

\paragraph{Fine-tuning the self-supervised ViT-B.}
Figure~\ref{fig:trackA-m7} shows that a vanilla ViT-B pre-trained
\emph{only on the challenge data} is viable but brittle. Layer-wise LR
decay with warm-up is the single most important ingredient: a naive
constant-LR fine-tune plateaus at $0.26$ held-out, whereas
LLRD$+$warm-up reaches $0.39$--$0.41$. Conversely, a poor pre-training
initialisation collapses fine-tuning to $0.027$ (near chance), so the
value of the self-supervised stage is entirely gated by its quality. The
persistent $\sim$17\,pt gap between the internal split and the true
held-out set mirrors the $\sim$28\,pt train/val gap of Track B and points
again to a data, not architecture, bottleneck.

\begin{figure}[t]
  \centering
  \safefig{fig_trackA_mae_pretrain.png}
  \caption{Self-supervised VideoMAE\,v2 pre-training from scratch
           (norm-pixel MSE). A $0.75$ tube-mask ratio converges smoothly
           to $\sim\!0.30$; pushing the encoder mask to $0.90$ without
           extra stabilisation diverges to NaN around epoch~29 (dotted).}
  \label{fig:trackA-mae}
\end{figure}

\paragraph{Pre-training stability.}
Figure~\ref{fig:trackA-mae} reports the masked-reconstruction loss. The
$0.75$ encoder mask ratio is stable and reaches a norm-pixel MSE of
$\sim\!0.30$, while the more aggressive $0.90$ ratio --- attractive for its
higher token-drop efficiency --- diverges to NaN early unless re-stabilised,
illustrating how sensitive from-scratch video MAE is at this data scale.

\begin{figure}[t]
  \centering
  \safefig{fig_trackA_confusion.png}
  \caption{Confusion matrix (row-normalised) of the best from-scratch
           model on \emph{val\_dir}. Red crosses mark the six most
           confused class pairs; as in Track B they are visually similar
           transitive actions separated only by a preposition or a
           direction, i.e.\ exactly where 4 frames carry too little motion
           information.}
  \label{fig:trackA-confusion}
\end{figure}

\paragraph{Error structure.}
The dominant confusions of the best from-scratch model
(Fig.~\ref{fig:trackA-confusion}) are semantic near-twins that 4 frames
cannot disambiguate: \emph{pretending} vs.\ real actions
(\emph{pretending to put something into something}
$\to$ \emph{putting something into something}, $40\%$;
\emph{pretending to throw} $\to$ \emph{throwing}, $23\%$),
preposition swaps (\emph{pouring something out of something}
$\to$ \emph{pouring something into something}, $23\%$) and direction
swaps (\emph{pulling right-to-left} $\to$ \emph{pulling left-to-right},
$23\%$). These are exactly the classes whose label lives in the motion
\emph{between} frames, which the sparse 4-frame sampling discards --- the
same failure mode observed for the pre-trained Track\,B models.

%==================================================================
\section{Conclusion and future work}
\label{sec:conclusion}

Our main findings are: \emph{(i)}~V-JEPA\,2 ViT-L (305\,M parameters)
is the best single-model architecture on Track\,B, reaching
$0.638$ \emph{val\_dir} accuracy and improving on the official
VideoMAE-B baseline ($0.6086$) by $+3.3$\,pts; \emph{(ii)}~scaling
beyond 300\,M parameters is unproductive at the data scale of
Track\,B, with 6--7\,B-parameter foundation models underperforming
both VideoMAE-B and V-JEPA\,2 ViT-L; \emph{(iii)}~on Track\,A, the best
from-scratch model tops out at $0.409$ --- some $20$\,pts below the
weakest pre-trained backbone --- and the R(2+1)D$+$Transformer matches the
self-supervised ViT-B at a fraction of its cost; \emph{(iv)}~the residual
generalisation gap ($\sim$28\,pts on Track\,B, $\sim$17\,pts on Track\,A)
cannot be closed by any of the five classes of regularisation we explored,
suggesting that more labelled data, not more architecture, is the binding
constraint.

Future directions include: cross-architecture ensembling with
per-class softmax weighting (where VideoMAE and V-JEPA-2 already
show complementary error patterns); self-supervised continued
pre-training on the unlabelled Track-B train partition; learning
the 4\,$\rightarrow$\,16 frame upsampling jointly with the
classifier instead of using a hand-crafted interpolation; and
distilling V-JEPA\,2 ViT-g into a smaller deployment-friendly head.

Another avenue we did not have time to explore
is to combine MGMAE~\cite{huang2023mgmae} with a strong optical-flow front-end such
as SEA-RAFT~\cite{wang2024searaft}. The 28-pt train/val gap we
report is partly a symptom of the model latching onto static
appearance cues --- texture, object identity, background --- rather
than the motion pattern that actually defines the
\emph{Something\,Something} classes (e.g.\ \texttt{Pulling left
to right} vs.\ \texttt{Pulling right to left}). Pre-computing
dense bidirectional flow between the four sampled frames with
SEA-RAFT (state-of-the-art accuracy on Sintel/KITTI at
$<10$\,ms per pair) and either (i) feeding the resulting flow
maps as an extra two-channel stream into a dual-pathway
VideoMAE, or (ii) using the SEA-RAFT flow as the motion prior
that drives MGMAE's motion-guided masking during continued
masked-pixel pre-training, would inject the explicit motion
signal that 4-frame sampling currently discards. We expect this
to shrink the train/val gap more effectively than any of the
appearance-based regularisers investigated in Sec.~\ref{sec:results},
and to disproportionately help the symmetric-pair classes that
dominate our confusion matrix.

%==================================================================
{\small
\begin{thebibliography}{99}

\bibitem{bai2023qwenvl}
J.~Bai, S.~Bai, S.~Yang, S.~Wang, S.~Tan, P.~Wang, J.~Lin, C.~Zhou, and J.~Zhou.
\newblock {Qwen-VL: A Versatile Vision-Language Model for Understanding, Localization, Text Reading, and Beyond}.
\newblock arXiv preprint arXiv:2308.12966, 2023.

\bibitem{bardes2024vjepa2}
A.~Bardes, Q.~Garrido, X.~Chen, M.~Assran, N.~Ballas, et~al.
\newblock {V-JEPA 2: Self-Supervised Video Models from One Million Hours of Internet Video}.
\newblock arXiv preprint, 2024.

\bibitem{bertasius2021timesformer}
G.~Bertasius, H.~Wang, and L.~Torresani.
\newblock {Is Space-Time Attention All You Need for Video Understanding?}
\newblock In \emph{International Conference on Machine Learning (ICML)}, 2021.

\bibitem{fan2021multiscale}
H.~Fan, B.~Xiong, K.~Mangalam, Y.~Li, Z.~Yan, J.~Malik, and C.~Feichtenhofer.
\newblock {Multiscale Vision Transformers}.
\newblock In \emph{IEEE/CVF International Conference on Computer Vision (ICCV)}, 2021.

\bibitem{goyal2017something}
R.~Goyal, S.~Ebrahimi~Kahou, V.~Michalski, J.~Materzynska, S.~Westphal, H.~Kim, et~al.
\newblock {The ``Something Something'' Video Database for Learning and Evaluating Visual Common Sense}.
\newblock In \emph{IEEE International Conference on Computer Vision (ICCV)}, 2017.

\bibitem{hu2022lora}
E.~J. Hu, Y.~Shen, P.~Wallis, Z.~Allen-Zhu, Y.~Li, S.~Wang, L.~Wang, and W.~Chen.
\newblock {LoRA: Low-Rank Adaptation of Large Language Models}.
\newblock In \emph{International Conference on Learning Representations (ICLR)}, 2022.

\bibitem{liu2023llava}
H.~Liu, C.~Li, Q.~Wu, and Y.~J. Lee.
\newblock {Visual Instruction Tuning}.
\newblock In \emph{Advances in Neural Information Processing Systems (NeurIPS)}, 2023.

\bibitem{liu2022videoswin}
Z.~Liu, J.~Ning, Y.~Cao, Y.~Wei, Z.~Zhang, S.~Lin, and H.~Hu.
\newblock {Video Swin Transformer}.
\newblock In \emph{IEEE/CVF Conference on Computer Vision and Pattern Recognition (CVPR)}, 2022.

\bibitem{ni2022xclip}
B.~Ni, H.~Peng, M.~Chen, S.~Zhang, G.~Meng, J.~Fu, S.~Xiang, and H.~Ling.
\newblock {Expanding Language-Image Pretrained Models for General Video Recognition}.
\newblock In \emph{European Conference on Computer Vision (ECCV)}, 2022.

\bibitem{oquab2023dinov2}
M.~Oquab, T.~Darcet, T.~Moutakanni, H.~Vo, M.~Szafraniec, V.~Khalidov, et~al.
\newblock {DINOv2: Learning Robust Visual Features without Supervision}.
\newblock arXiv preprint arXiv:2304.07193, 2023.

\bibitem{ryali2023hiera}
C.~Ryali, Y.-T. Hu, D.~Bolya, C.~Wei, H.~Fan, P.-Y. Huang, V.~Aggarwal, A.~Chowdhury, O.~Poursaeed, J.~Hoffman, J.~Malik, Y.~Li, and C.~Feichtenhofer.
\newblock {Hiera: A Hierarchical Vision Transformer without the Bells-and-Whistles}.
\newblock In \emph{International Conference on Machine Learning (ICML)}, 2023.

\bibitem{szegedy2016rethinking}
C.~Szegedy, V.~Vanhoucke, S.~Ioffe, J.~Shlens, and Z.~Wojna.
\newblock {Rethinking the Inception Architecture for Computer Vision}.
\newblock In \emph{IEEE Conference on Computer Vision and Pattern Recognition (CVPR)}, 2016.

\bibitem{tong2022videomae}
Z.~Tong, Y.~Song, J.~Wang, and L.~Wang.
\newblock {VideoMAE: Masked Autoencoders are Data-Efficient Learners for Self-Supervised Video Pre-Training}.
\newblock In \emph{Advances in Neural Information Processing Systems (NeurIPS)}, 2022.

\bibitem{tran2018closer}
D.~Tran, H.~Wang, L.~Torresani, J.~Ray, Y.~LeCun, and M.~Paluri.
\newblock {A Closer Look at Spatiotemporal Convolutions for Action Recognition}.
\newblock In \emph{IEEE/CVF Conference on Computer Vision and Pattern Recognition (CVPR)}, 2018.

\bibitem{lin2019tsm}
J.~Lin, C.~Gan, and S.~Han.
\newblock {TSM: Temporal Shift Module for Efficient Video Understanding}.
\newblock In \emph{IEEE/CVF International Conference on Computer Vision (ICCV)}, 2019.

\bibitem{tschannen2025siglip2}
M.~Tschannen, A.~Gritsenko, X.~Wang, M.~F. Naeem, I.~Alabdulmohsin, et~al.
\newblock {SigLIP 2: Multilingual Vision-Language Encoders with Improved Semantic Understanding, Localization, and Dense Features}.
\newblock arXiv preprint arXiv:2502.14786, 2025.

\bibitem{wang2024internvideo2}
Y.~Wang, K.~Li, X.~Li, J.~Yu, Y.~He, G.~Chen, B.~Pei, R.~Zheng, et~al.
\newblock {InternVideo2: Scaling Foundation Models for Multimodal Video Understanding}.
\newblock In \emph{European Conference on Computer Vision (ECCV)}, 2024.

\bibitem{zhang2018mixup}
H.~Zhang, M.~Cisse, Y.~N. Dauphin, and D.~Lopez-Paz.
\newblock {mixup: Beyond Empirical Risk Minimization}.
\newblock In \emph{International Conference on Learning Representations (ICLR)}, 2018.

\bibitem{wang2024searaft}
Y.~Wang, L.~Lipson, and J.~Deng.
\newblock {SEA-RAFT: Simple, Efficient, Accurate RAFT for Optical Flow}.
\newblock In \emph{European Conference on Computer Vision (ECCV)}, 2024.

\bibitem{huang2023mgmae}
B.~Huang, Z.~Zhao, G.~Zhang, Y.~Qiao, and L.~Wang.
\newblock {MGMAE: Motion Guided Masking for Video Masked Autoencoding}.
\newblock In \emph{IEEE/CVF International Conference on Computer Vision (ICCV)}, 2023.

\end{thebibliography}
}

\appendix
\section{Classical data augmentation}
\label{sec:augmentation}

To narrow the train/val gap reported in
Sec.~\ref{sec:method:overfit}, every training run applies a stack of
classical, label-preserving image transformations to each of the four
source frames \emph{before} the upsampling step described in
Sec.~\ref{sec:upsampling}. The same random parameters are sampled
\emph{once per clip} and broadcast across the four frames, so that
temporal coherence — the very signal that distinguishes \emph{Pulling
left to right} from \emph{Pulling right to left} — is preserved.

\begin{figure*}[!ht]
  \centering
  \safefig[\textwidth]{fig_augmentation_real.png}
  \caption{Seven classical data-augmentation transforms applied to one
           frame of the Track\,B video shown in
           Fig.~\ref{fig:upsampling} (class~3, \emph{Folding
           something}). The same random parameters are shared across
           the four source frames of a clip so that motion cues are
           not destroyed. \emph{Colour jitter} and \emph{random
           resized crop} are the largest individual contributors
           ($+1.4$ and $+1.1$\,pt val\_dir respectively), MixUp adds a
           further $+0.6$\,pt, and temporal-reverse doubles the
           effective sample count for the 16 directional classes.}
  \label{fig:augmentation}
\end{figure*}

\section{Frame upsampling strategies}
\label{sec:upsampling}

The VideoMAE and V-JEPA\,2 tube tokeniser (kernel 2 frames~$\times$~16$\times$16~px)
expects 16 input frames per clip, while Track\,B delivers only 4. We
compared six strategies to bridge this gap before feeding the
backbone (Fig.~\ref{fig:upsampling}).

\begin{figure*}[!ht]
  \centering
  \safefig[\textwidth]{fig_frame_upsampling_real.png}
  \caption{Six strategies for upsampling the 4~source frames to the
           sequence expected by VideoMAE-B and V-JEPA\,2's tube tokeniser.
           Each output position is colour-coded by the convex combination
           of the four source colours dictated by the method. The methods' efficiency on VideoMAE-B top1 accuracy are sumarized in Table \ref{tab:upsampling} }
  \label{fig:upsampling}
\end{figure*}

\begin{table}[!ht]
  \centering
  \small
  \begin{tabular}{@{}lcccccc@{}}
    \toprule
    method
                    & replicate         & nearest & cubic & Gaussian & LERP & tile  \\
    \midrule
    \emph{val\_dir} & \textbf{0.6086}   & 0.598   & 0.588 & 0.582 & 0.54     & 0.412 \\
    \bottomrule
  \end{tabular}
  \caption{Top-1 \emph{val\_dir} accuracy for each upsampling strategy.
           The $-6.6$\,pt drop between replicate and LERP is consistent
           with the SSV2 distribution being dominated by sharp natural
           frames rather than blended intermediates.}
  \label{tab:upsampling}
\end{table}

\end{document}
